\documentclass[12pt,a4paper]{article}
\usepackage[spanish]{babel}
\usepackage[margin=2.5cm]{geometry}
\usepackage{graphicx}
\usepackage{listings}
\usepackage{xcolor}
\usepackage{amsmath}
\usepackage[hidelinks]{hyperref}

\definecolor{codebg}{RGB}{245,245,245}

\lstset{
    language=Python,
    basicstyle=\ttfamily\small,
    backgroundcolor=\color{codebg},
    frame=single,
    rulecolor=\color{codebg},
    framesep=5pt,
    keywordstyle=\color{blue},
    commentstyle=\color{gray},
    stringstyle=\color{red},
    breaklines=true,
    showstringspaces=false
}

\title{Optimización de un Autocorrector \\ Memoria Técnica}
\author{Eloy Urriens - Juan Pulgarín - Tomas Alignani - Nicolás Chareca}
\date{22 de marzo de 2026}

\begin{document}

\maketitle

\begin{abstract}
Presentamos el proceso de desarrollo y optimización de un autocorrector basado inicialmente en unigramas y distancia de edición. A partir de esa implementación base, analizamos sus limitaciones y planteamos mejoras orientadas a incorporar información contextual mediante n-gramas y a refinar el coste de las operaciones de edición. El objetivo final es obtener un modelo más preciso y robusto, evaluado sobre conjuntos de validación y test independientes.
\end{abstract}

\clearpage

\tableofcontents
\newpage

\section{Introducción}

Este proyecto de la asignatura 'Sistemas Inteligentes. Aplicaciones' tiene como objetivo la optimización de un modelo de autocorrección ortográfica. Partimos de una implementación básica que permite la detección de palabras fuera de vocabulario, utilizando la distancia de Levenshtein para la selección de la corrección más probable.

La implementación base, desarrollada para la práctica, utiliza un corpus de entrenamiento compuesto por 21 libros de Pío Baroja para calcular las probabilidades de aparición de cada palabra. Sin embargo, este modelo presenta limitaciones como la falta de sensibilidad al contexto y la incapacidad de detectar errores en los que la palabra escrita existe en el diccionario pero es incorrecta en su entorno (ej: la gran 'automóvil').

El propósito de este trabajo la implementación de una serie de mejoras para este corrector que lo transformen a un modelo más sofisticado. Para ello, planteamos las siguientes mejoras:

\begin{itemize}
    \item \textbf{\underline{Contexto} / n-gramas:} Evolucionar el modelo de unigramas hacia bigramas o trigramas para capturar información contextual y mejorar la precisión de las sugerencias.  permitirá corregir errores donde la palabra escrita existe en el diccionario pero no tiene sentido en la oración.
    \item \textbf{\underline{Distancia} de edición ponderada por teclado:} Modificar la penalización de las ediciones para considerar la cercanía física de las teclas en un teclado QWERTY.
\end{itemize}

Para evaluar el impacto de cada mejora, definimos conjuntos de datos de validación y test generales para usarse por los métodos mencionados, permitiéndonos experimentar con los hiperparámetros del mismo (ej: n de los ngramas) y así obtener datos reales del funcionamiento de estos modelos con distintas configuraciones, para lograr un mejor entendimiento de los mismos.

\section{Análisis del Sistema Actual}
\subsection{Funcionamiento Base}

\paragraph{Construcción del modelo de lenguaje.}
Como se mencionó en la introducción, el corpus se forma por 21 libros de Pío Baroja. El texto se normaliza a minúsculas y se eliminan todos los caracteres no alfanuméricos con \lstinline[language={}]|re.sub(r"[^\w\s]", "", texto)|. La probabilidad de cada palabra se calcula como su frecuencia relativa sobre el tamaño total del corpus:

\begin{equation}
    P(w) = \frac{\text{C}(w)}{\text{V}}
\end{equation}

\begin{itemize}
    \item $P(w)$: probabilidad de la palabra $w$.
    \item $C(w)$: número de veces que aparece la palabra $w$ en el corpus.
    \item $V$: tamaño total del corpus.
\end{itemize}

\paragraph{Generación de palabras candidatas a X distancia de edición.}
Dada una palabra de entrada, se generan todos los strings obtenibles (pertenecientes o no al corpus) aplicando las cuatro operaciones elementales:

\begin{lstlisting}
def una_edicion(word: str, intercambiar: bool=False) -> set:
    words_insert = insertar(word)
    words_remove = borrar(word)
    words_replace = reemplazar(word)
    words_swap = intercambiar(word) if intercambiar else []
    words = set(words_insert + words_remove + words_replace + words_swap)
    return words
\end{lstlisting}

\begin{lstlisting}
def dos_ediciones(word: str, intercambiar: bool=False) -> set:
    words_dos_ediciones = set()
    words_una_edicion = una_edicion(word, intercambiar)
    for word_una_ed in words_una_edicion:
        words_dos_ediciones.union(una_edicion(word_una_ed, intercambiar))
    return words_dos_ediciones
\end{lstlisting}

La función \texttt{corregir()} filtra las palabras generadas conservando 
solo las que están presentes en el corpus, luego las ordena por $P(w)$ 
(probabilidad de aparición) de mayor a menor y devuelve las $n$ más probables. 
Priorizamos la distancia 1 sobre la 2, ya que es más probable que 
se equivoque 1 vez a 2 veces, y si la palabra original ya 
pertenece al corpus, se retorna directamente sin buscar correcciones.

\paragraph{Distancia de edición mínima entre frases.}
Para comparar la frase original con la corregida implementamos el algoritmo de \href{https://en.wikipedia.org/wiki/Wagner\%E2\%80\%93Fischer_algorithm}{Wagner-Fischer} mediante programación dinámica. La matriz de costes $M$ de dimensiones ($size_{frase1}+1 \times size_{frase2}+1$), en su posicion $[i, j]$ se rellena con la distancia de Levenshtein de la frase1[:i] y frase2[:j], utilizando la siguiente fórmula:

\begin{equation}
    M[i,j] = \begin{cases}
        M[i-1,\,j-1] & \text{si } s_1[i] = s_2[j] \\[4pt]
        1 + \min\bigl(M[i-1,j],\; M[i,j-1],\; M[i-1,j-1]\bigr) & \text{en otro caso}
    \end{cases}
\end{equation}

El coste total de corrección de la frase es la suma de los costes individuales de cada par de palabras, obtenible en la posición M[$size_{frase1}, size_{frase2}$]. 

Luego por último, implementamos una función recursiva encargada de obtener los pasos necesarios para transformar la frase original a la corregida, utilizando la matriz de costes $M$ y aplicando backtracking desde la posición M[$size_{frase1}, size_{frase2}$] hasta M[0, 0].

\clearpage

\section{Propuestas de Mejora}
\subsection{Algoritmos Implementados}
\subsubsection{Contexto / n-gramas}

\paragraph{Introducción.}
Para esta mejora al modelo de autocorrección, se propone evolucionar el modelo de unigramas hacia bigramas o trigramas. Esto permitirá capturar información contextual y mejorar la precisión de las sugerencias, especialmente en casos donde la palabra escrita existe en el diccionario pero no tiene sentido en la oración (ej: "la gran 'automóvil'").

\paragraph{Implementación.}
El algoritmo que implementamos sigue los siguientes pasos:

\begin{enumerate}
    \item Llega una frase a corregir
    \item Para cada palabra de la frase
        \subitem Si esta en el vocabulario --> La dejamos como esta
        \subitem Si no está en el vocabulario --> Obtenemos todas las palabras cercanas y frecuentes junto con sus probabilidades de contexto (probabilidades que indican como de bien encajan esas palabras allí) y luego se elige la de mayor probabilidad.
    \item Devolvemos la frase corregida
\end{enumerate}

\paragraph{Hiperparámetros.}
Esta mejora utiliza 2 hiperparámetros que son \textbf{min\_appearances} (mínima cantidad de apariciones de la frase en el corpus para no ser reemplazada por el token \texttt{<UNK>}) y \textbf{n} (número de palabras por n-grama).

\paragraph{Validación de hiperparámetros.}
Para validar los hiperparámetros, definimos un conjunto de validación específico, compuesto por frases con errores ortográficos. Se experimenta con el rango $[1;20]$ (step=1) de \textbf{min\_appearances} para evaluar su impacto en la precisión de las correcciones, utilizando la metrica de accuracy\_score, donde si la frase corregida es igual a la frase correcta, entonces suma 1. \textbf{n} se mantiene como 2, ya que para n mayor se requiere una alta cantidad de memoria para almacenar las matrices.



\subsubsection{Distancia física de teclado}

\subsection{Estructura del Código}
\begin{lstlisting}
# Poner codigo
\end{lstlisting}

\section{Resultados}
Presentar métricas de desempeño y comparativas.

\subsection{Resultados por método}
\subsubsection{Contexto / n-gramas}

\paragraph{Resultados de la validación realizada.}
Encontramos que para \textbf{min\_appearances} igual a 2, 3, 4, 6 y 11 se obtiene el mejor resultado (\texttt{accuracy\_score} 40\%), lo que indica que el modelo corrige correctamente el 40\% de las frases del conjunto de validación.

Esto señala que el modelo requiere un balance entre capturar patrones frecuentes del corpus y evitar ser demasiado restrictivo. 

Un \texttt{min\_appearances} bajo (como 1-2) permite corregir más casos, pero puede introducir ruido de n-gramas poco representativos. Un valor más alto (como 11) filtra mejor el ruido pero puede perder correcciones válidas. 

Nos quedamos con \texttt{min\_appearances} = 11 para la evaluación final, ya que es el valor más alto que mantiene el mejor resultado, lo cual nos permite obtener un modelo con una mejor capacidad de generalización, y más flexibilidad ante textos con palabras nuevas.

Este resultado sugiere que la mejora con n-gramas es moderada, probablemente porque:
\begin{enumerate}
    \item El corpus de entrenamiento es limitado, lo que restringe la cobertura de bigramas.
    \item El contexto por sí solo no es suficiente sin combinar otras señales como la distancia de edición ponderada.
\end{enumerate} 

Aún contando con un corpus mayor, los bigramas no aportan suficiente información contextual. Por ejemplo, frases que comiencen por 'el': el pajaro, el raton, el coche, el ordenador, son bigramas que no aportan información útil para el contexto.

\subsubsection{Distancia física de teclado}

\subsubsection{Combinacion (Contexto + Distancia)}

\subsection{Tabla comparativa}


\section{Conclusiones}
Resumen de mejoras logradas y trabajo futuro.

\end{document}